%% Correction audit de veracite - 2026-09-20 : citations de rapport de jury sans
%% support textuel dans le corpus (rapports 2018-2025) remplacees par des verbatims
%% verifies ou supprimees. Detail ci-dessous.
%% l.312 (ex) : pseudo-citation "Rapport jury 2024 : ne pas omettre la partie reactive"
%%   INTROUVABLE dans tout le corpus -> supprimee (le constat technique est conserve).
%% Les 3 citations de l'encadre jury (RJ2024 / RJ2021 / RJ2020) sont conformes : non touchees.
\expandafter\ifx\csname CAPETINCLUDEMODE\endcsname\relax
\documentclass[11pt,a4paper]{article}
% preamble.tex — packages et commandes communes a toutes les fiches CAPET SII IE
\input{glyphtounicode}
\pdfgentounicode=1
\usepackage{cmap}
\usepackage[T1]{fontenc}
\usepackage[utf8]{inputenc}
\usepackage[french]{babel}
% babel[french] active a \begin{document} des raccourcis de ponctuation (:;!?)
% qui cassent la syntaxe to[...] de circuitikz (erreur "+ or - expected") ;
% on les desactive, ce qui ne change que l'espacement fin avant ces signes.
\AtBeginDocument{\shorthandoff{:;!?}}
\usepackage[a4paper, top=1.6cm, bottom=1.6cm, left=1.4cm, right=1.4cm]{geometry}
\setlength{\headheight}{14.5pt}
\usepackage{amsmath, amssymb, amsfonts}
\usepackage{siunitx}
\sisetup{locale=FR, output-decimal-marker={,}, inter-unit-product=\ensuremath{{}\cdot{}},
         group-minimum-digits=4}
% Unites non fournies par siunitx, declarees ici une fois pour toutes (fiches + exercices).
% Les gardes \ifdefined evitent tout conflit avec une redeclaration locale restee dans un
% fichier : ne pas retirer les gardes, ne pas redeclarer ces unites dans les fichiers.
\DeclareSIUnit{\jour}{j}
\ifdefined\var\else\DeclareSIUnit{\var}{var}\fi   % puissance reactive
\ifdefined\vad\else\DeclareSIUnit{\vad}{VAD}\fi   % puissance deformante
\ifdefined\tr\else\DeclareSIUnit{\tr}{tr}\fi      % tour (tr/min)
\usepackage{array, booktabs, multirow, tabularx}
\usepackage{xcolor}
\usepackage[most, breakable]{tcolorbox}
\usepackage{multicol}
\usepackage{enumitem}
\usepackage{mathtools}
\usepackage{lmodern}
% Refonte visuelle 2026 (direction "moderne coloré") : Lato pour le texte,
% maths sans-serif assorties via newtxsf (charge apres amsmath/amssymb).
\usepackage[default]{lato}
\usepackage{newtxsf}
\usepackage{microtype}
\usepackage{fancyhdr}
\usepackage{tikz}
\usetikzlibrary{positioning,arrows.meta}
\usepackage{circuitikz}
\usepackage{lastpage}

% ── Couleurs par domaine (palette modernisée 2026) ────────────────────
\definecolor{colElec}{RGB}{29, 78, 216}        % Électrotechnique — bleu vif
\definecolor{colElecLight}{RGB}{219, 234, 254} % fond bleu clair
\definecolor{colAuto}{RGB}{21, 128, 61}        % Automatique — vert
\definecolor{colMeca}{RGB}{194, 65, 12}        % Mécanique — orange
\definecolor{colTherm}{RGB}{185, 28, 28}       % Thermique — rouge
\definecolor{colQT}{RGB}{232, 244, 233}        % (hérité, plus utilisé par questiontype)
\definecolor{colJury}{RGB}{255, 248, 225}      % fond jury (crème)
\definecolor{colJuryFrame}{RGB}{202, 138, 4}   % bordure/bandeau jury (ambre)
\definecolor{colGrayLight}{RGB}{245, 245, 247} % fond gris neutre
\definecolor{colSeparator}{RGB}{180, 180, 180} % gris séparateur
\definecolor{colExo}{RGB}{8, 145, 178}          % Banque d'exercices — cyan (distinct des fiches)
% ── Encre et gris de la refonte ───────────────────────────────────────
\definecolor{colInk}{RGB}{23, 32, 42}          % texte courant (quasi-noir)
\definecolor{colSlate}{RGB}{83, 98, 118}       % gris ardoise (méta, en-têtes)
\definecolor{colRule}{RGB}{203, 213, 225}      % filets clairs
\definecolor{colJuryBg}{RGB}{255, 248, 225}    % fond crème jury
\definecolor{colPresence}{RGB}{83, 98, 118}    % bloc 5 (présence dans les sujets)
\AtBeginDocument{\color{colInk}}

% ── Style encadré exercice + corrigé (banque d'exercices) ─────────────
\tcbset{
  exoBox/.style={
    breakable,
    enhanced,
    colback=colExo!8!white,
    colframe=colExo,
    colbacktitle=colExo,
    coltitle=white,
    fonttitle=\bfseries\small,
    left=4pt, right=4pt, top=3pt, bottom=3pt,
    before skip=5pt, after skip=5pt,
  }
}

% ── Bascule enonce / corrige (banque d'exercices) ─────────────────────
% \CAPETCORRIGE est defini en ligne de commande (pdflatex "\def\CAPETCORRIGE{1}\input{...}")
% pour compiler la variante corriges d'un meme fichier source ; par defaut (non defini),
% la compilation produit la variante enonces. \CAPETSHOW{enonce}{corrige} choisit la branche.
%
% TROIS MODES (refonte 2026-09-20, manuel de revision) :
%   - aucun drapeau   -> imprime l'enonce            (variante enonces, INCHANGE)
%   - \CAPETCORRIGE   -> imprime le corrige          (variante corriges, INCHANGE)
%   - \CAPETMANUEL    -> imprime l'enonce ET met le corrige de cote pour l'annexe
%
% Le mode manuel stocke chaque corrige dans une macro (\csgdef), sans ecriture
% fichier ni re-tokenisation : le contenu (circuitikz, accents, %) traverse intact.
% L'annexe les rejoue ensuite dans l'ordre via \CAPETANNEXE.
\usepackage{etoolbox}

\newcounter{CAPETexo}          % numero continu d'exercice dans le manuel
\setcounter{CAPETexo}{0}
\newcounter{CAPETannexe}       % compteur de relecture en annexe

% Titre courant de la notion en cours (pose par \CAPETNOTION), sert d'intertitre annexe
\newcommand{\CAPETnotioncode}{}
\newcommand{\CAPETnotiontitre}{}

% \CAPETNOTION{code}{titre} : ouvre une notion dans le manuel (memorise le libelle,
% et marque le rang du premier exercice de la notion pour l'intertitre d'annexe).
\newcommand{\CAPETNOTION}[2]{%
  \renewcommand{\CAPETnotioncode}{#1}%
  \renewcommand{\CAPETnotiontitre}{#2}%
  \csgdef{CAPETnotionstart\the\numexpr\value{CAPETexo}+1\relax}{#1 -- #2}%
}

% ── \degree robuste ───────────────────────────────────────────────────
% Certains fichiers d'exercices utilisent \degree. Ni le preambule ni siunitx
% ne le definissent, et \providecommand local ne suffit pas : sous hyperref,
% \degree existe deja comme commande factice (\ERROR), donc \providecommand
% s'abstient et l'appel echoue a la compilation du manuel (erreur invisible
% en compilation isolee). On le (re)definit donc inconditionnellement ici,
% apres le chargement des packages, pour tout le corpus.
\AtBeginDocument{\providecommand{\degree}{}\renewcommand{\degree}{\ensuremath{^\circ}}}

% Libelle decoratif : les fichiers d'exercices utilisent aussi \CAPETSHOW dans
% leur \fancyhead et leur bloc de titre, pour afficher "(enonces)" ou "(corriges)".
% Ces appels-la ne sont PAS des exercices : en mode manuel ils ne doivent ni
% consommer un numero, ni alimenter l'annexe. On les reconnait a leur premier
% argument, strictement egal au libelle ci-dessous.
\newcommand{\CAPETLIBELLEENONCES}{(\'enonc\'es)}

\ifdefined\CAPETMANUEL
  % --- Mode manuel : enonce en place, corrige differe vers l'annexe ---
  \newcommand{\CAPETSHOW}[2]{%
    \def\CAPETargun{#1}%
    \ifx\CAPETargun\CAPETLIBELLEENONCES
      #1%            appel decoratif : rendu tel quel, hors numerotation
    \else
      \CAPETSHOWEXO{#1}{#2}%
    \fi
  }
  \newcommand{\CAPETSHOWEXO}[2]{%
    \stepcounter{CAPETexo}%
    \csgdef{CAPETcorrige\theCAPETexo}{#2}%
    \edef\CAPETtmpnum{\theCAPETexo}%
    #1%
    \par\vspace{2pt}%
    {\footnotesize\itshape\color{colExo}%
     \hyperref[capet:cor:\CAPETtmpnum]{Corrige de l'exercice \CAPETtmpnum\ $\rightarrow$ p.~\pageref{capet:cor:\CAPETtmpnum}}%
     \label{capet:eno:\CAPETtmpnum}}%
  }
  % \CAPETANNEXE : rejoue tous les corriges captures, dans l'ordre, avec intertitres
  \newcommand{\CAPETANNEXE}{%
    \setcounter{CAPETannexe}{0}%
    \loop\ifnum\value{CAPETannexe}<\value{CAPETexo}%
      \stepcounter{CAPETannexe}%
      \edef\CAPETcurnum{\theCAPETannexe}%
      \ifcsname CAPETnotionstart\CAPETcurnum\endcsname
        \par\medskip
        {\normalsize\bfseries\color{colExo}\csuse{CAPETnotionstart\CAPETcurnum}}%
        \par\nobreak\smallskip
      \fi
      \begin{tcolorbox}[exoBox, breakable,
        title={Corrige de l'exercice \CAPETcurnum\
               \hfill\normalfont\footnotesize
               \hyperref[capet:eno:\CAPETcurnum]{enonce p.~\pageref{capet:eno:\CAPETcurnum}}}]%
      \label{capet:cor:\CAPETcurnum}%
      \csuse{CAPETcorrige\CAPETcurnum}%
      \end{tcolorbox}%
    \repeat
  }
\else
  % --- Modes historiques, comportement strictement inchange ---
  \newcommand{\CAPETSHOW}[2]{\ifdefined\CAPETCORRIGE #2\else #1\fi}
  \newcommand{\CAPETANNEXE}{}
\fi

% ── Difficulté 1 à 3 (étoiles pleines, amssymb) ───────────────────────
\newcommand{\difficulte}[1]{%
  \ifcase#1\relax
  \or $\bigstar$\;niveau 1/3%
  \or $\bigstar\bigstar$\;niveau 2/3%
  \or $\bigstar\bigstar\bigstar$\;niveau 3/3%
  \fi
}

% ── Style encadré extrait de document technique (banque d'exercices) ──
\tcbset{
  docTech/.style={
    breakable,
    enhanced,
    colback=colGrayLight,
    colframe=colSeparator,
    fonttitle=\bfseries\small\itshape,
    coltitle=black,
    colbacktitle=colGrayLight,
    left=4pt, right=4pt, top=3pt, bottom=3pt,
    before skip=4pt, after skip=4pt,
  }
}

% ── Style tcolorbox de section principale ────────────────────────────
% Refonte 2026 : bandeau titre coloré arrondi + fond très pâle, sans cadre.
\tcbset{
  ficheSection/.style={
    breakable, enhanced jigsaw,
    frame hidden, boxrule=0pt, arc=1.8mm,
    colback=#1!4!white,
    colbacktitle=#1,
    coltitle=white,
    fonttitle=\sffamily\bfseries,
    left=7pt, right=7pt, top=5pt, bottom=5pt,
    toptitle=3pt, bottomtitle=3pt,
    before skip=10pt, after skip=8pt,
  }
}

% ── Style encadré question type ──────────────────────────────────────
% Carte blanche, bandeau sombre, liseré très clair.
\tcbset{
  questiontype/.style={
    breakable, enhanced jigsaw,
    colback=white,
    frame hidden, boxrule=0pt, arc=1.8mm,
    colbacktitle=colElec!78!black,
    coltitle=white,
    fonttitle=\sffamily\bfseries\small,
    left=7pt, right=7pt, top=4pt, bottom=4pt,
    toptitle=2.5pt, bottomtitle=2.5pt,
    before skip=8pt, after skip=8pt,
    borderline={0.6pt}{0pt}{colElec!12!white},
  }
}

% ── Style encadré jury ───────────────────────────────────────────────
% Bandeau ambre, fond crème, sans cadre.
\tcbset{
  jury/.style={
    breakable, enhanced jigsaw,
    colback=colJuryBg,
    frame hidden, boxrule=0pt, arc=1.8mm,
    colbacktitle=colJuryFrame,
    coltitle=white,
    fonttitle=\sffamily\bfseries,
    left=7pt, right=7pt, top=4pt, bottom=4pt,
    toptitle=3pt, bottomtitle=3pt,
    before skip=10pt, after skip=8pt,
  }
}

% ── En-tête / pied de page communs (refonte 2026) ─────────────────────
% #1 code fiche, #2 titre court, #3 couleur domaine
\newcommand{\ficheHeaders}[3]{%
  \fancyhead[L]{\small\sffamily{\bfseries\textcolor{#3}{#1}}\hspace{0.7em}\textcolor{colSlate}{#2}}%
  \fancyhead[R]{\small\sffamily\textcolor{colSlate}{CAPET SII IE\, --- \,Session 2026}}%
  \fancyfoot[C]{\small\sffamily\textcolor{colSlate}{\thepage\,/\,\pageref{LastPage}}}%
  \fancyhead[C]{}\fancyfoot[L]{}\fancyfoot[R]{}%
}

% ── Bloc en-tête de fiche (BLOC 0, refonte 2026) ──────────────────────
% #1 code  #2 titre  #3 couleur domaine  #4 domaine  #5 ligne présence  #6 sous-thèmes
\newcommand{\ficheEntete}[6]{%
  \begin{tcolorbox}[enhanced, breakable,
    frame hidden, boxrule=0pt, arc=2.5mm,
    interior style={left color=#3, right color=#3!75!black},
    left=9pt, right=9pt, top=7pt, bottom=7pt,
    before skip=0pt, after skip=6pt]
    {\color{white}\sffamily\bfseries\LARGE #1\quad #2}\\[4pt]
    {\color{#3!22!white}\sffamily\small \MakeUppercase{#4}\quad\textbar\quad #5}\\[1pt]
    {\color{#3!35!white}\sffamily\footnotesize Sous-thèmes : #6}
  \end{tcolorbox}%
}

% ── Mise en page ─────────────────────────────────────────────────────
\pagestyle{fancy}
\fancyhf{}
\renewcommand{\headrulewidth}{0.5pt}
\renewcommand{\footrulewidth}{0pt}
\AtBeginDocument{%
  \renewcommand{\headrule}{{\color{colRule}\hrule width\headwidth height 0.5pt}}%
}

\setlist[itemize]{leftmargin=1.2em, itemsep=1pt, topsep=2pt, parsep=0pt}
\setlist[enumerate]{leftmargin=1.2em, itemsep=3pt, topsep=2pt, parsep=0pt}

\setlength{\parskip}{2pt}
\setlength{\columnsep}{10pt}

\def\CAPETBEGINFILE{\begin{document}}
\def\CAPETENDFILE{\end{document}}
\else
\def\CAPETBEGINFILE{}
\def\CAPETENDFILE{}
\fi

\ficheHeaders{T17}{Distribution et sécurité}{colElec}

\CAPETBEGINFILE

% -----------------------------------------------------------------------
%  BLOC EN-TETE
% -----------------------------------------------------------------------

\ficheEntete{T17}{Distribution et sécurité électrique, régimes de neutre}{colElec}%
  {électrotechnique - Distribution BT}%
  {\textbf{Présence :} \textbf{5/13} sujets (2013-2026) - thème \emph{régulier}}%
  {SLT (TT / TN-S / TN-C / IT) - DDR et sélectivité différentielle - disjoncteur MA devant variateur - chute de tension NF C 15-100 - parafoudres types 1/2/3 - bilan de puissance et abonnement}

% -----------------------------------------------------------------------
%  1. RAPPELS DE COURS
% -----------------------------------------------------------------------

\begin{tcolorbox}[
  ficheSection=colElec,
  title={1. Rappels de cours},
]

% -- 1.1 Schemas de liaison a la terre --------------------------------
\textbf{1.1 Schémas de liaison à la terre (SLT) - NF C 15-100}

\smallskip
\begin{center}
\footnotesize
\renewcommand{\arraystretch}{1.3}
\begin{tabular}{@{}p{1.1cm}p{1.9cm}p{2.6cm}p{3.3cm}p{2.9cm}@{}}
\toprule
\textbf{Régime} & \textbf{Neutre source} & \textbf{Masses récepteurs} & \textbf{Protection contacts indirects} & \textbf{Application} \\
\midrule
TT & Terre directe & Terre indépendante & \textbf{DDR obligatoire} & Résidentiel, tertiaire \\
TN-S & Terre directe & Reliées au neutre (PE séparé) & Disjoncteur ($I_{cc}$ élevé) & Industriel, scolaire \\
TN-C & Terre directe & Reliées au neutre (PEN) & Disjoncteur (DDR \emph{interdit}) & Anciens réseaux \\
IT & Impédant & Terre & CPI + coupure au 2\textsuperscript{e} défaut & Médical, industrie critique \\
\bottomrule
\end{tabular}
\normalsize
\end{center}

\medskip\hrule\vspace{4pt}

% -- 1.2 DDR et selectivite -------------------------------------------
\textbf{1.2 Dispositifs différentiels résiduels (DDR) - sensibilités normalisées}

\smallskip
\begin{center}
\begin{tabular}{lll}
\toprule
\textbf{Sensibilité} & \textbf{Usage} & \textbf{Norme} \\
\midrule
30\,mA (haute sensibilité) & Protection des personnes & NF C 15-100 - obligatoire en TT \\
300\,mA ou 500\,mA & Protection incendie, sélectivité & - \\
1\,A et plus & Protection de tête (sélectivité amont) & - \\
\bottomrule
\end{tabular}
\end{center}
\textbf{Sélectivité différentielle :} rapport des sensibilités DDR amont/aval $\geq 3$ (ex : DDR 30\,mA aval + DDR 300\,mA amont ou DDR temporisé de type S en amont).

\medskip\hrule\vspace{4pt}

% -- 1.3 Chute de tension ---------------------------------------------
\textbf{1.3 Chute de tension - NF C 15-100}

\smallskip
\[
  \Delta U\,\% = \frac{\sqrt{3}\,I\,L\,(R_u\cos\varphi + X_u\sin\varphi)}{U_D} \times 100
\]

\smallskip
\begin{center}
\begin{circuitikz}[european, scale=0.9]
  \draw (0,0) node[left]{$U_D$ (source)} to[short] (0.6,0)
        to[R, l=$R_uL$] (3,0)
        to[cute inductor, l=$X_uL$] (5.4,0)
        to[short] (6,0) node[right]{$U$ (récepteur)};
  \draw[<->] (0.6,-0.9) -- (5.4,-0.9);
  \node[below] at (3,-0.9) {\scriptsize $\Delta U = U_D - U$};
\end{circuitikz}
\end{center}

Valeurs limites : domaine public éclairage = \textbf{3\,\%} ; domaine public autres usages = \textbf{5\,\%}.\\
Si $\Delta U > 5\,\%$ : \textbf{augmenter la section} du câble pour réduire $R_u = \rho/S$.

\medskip\hrule\vspace{4pt}

% -- 1.4 Disjoncteur amont variateur ----------------------------------
\textbf{1.4 Disjoncteur magnétique (MA) en amont d'un variateur}

\smallskip
\begin{center}
\begin{tabular}{lll}
\toprule
\textbf{Critère} & \textbf{Valeur/Type} & \textbf{Justification} \\
\midrule
Type de déclencheur & MA (magnétique seul) & Variateur assure la protection thermique \\
Seuil magnétique $I_{rm}$ & Ex : 2860\,A (NSX250N MA220) & Court-circuit uniquement \\
Pouvoir de coupure $I_{cu}$ & 50\,kA & Courant de court-circuit réseau \\
\bottomrule
\end{tabular}
\end{center}
Un disjoncteur magnéto-thermique (MA+thermique) en amont créerait une double protection thermique - risque de décoordination avec le variateur.

\end{tcolorbox}

% -----------------------------------------------------------------------
%  2. FORMULAIRE SYNTHETIQUE
% -----------------------------------------------------------------------

\begin{tcolorbox}[
  ficheSection=colElec,
  title={2. Formulaire synthétique},
]
\small
\begin{multicols}{2}
\renewcommand{\arraystretch}{1.45}
\begin{tabular}{@{}p{4.5cm}ll@{}}
\toprule
\textbf{Formule} & \textbf{Usage} & \textbf{U.} \\
\midrule
\multicolumn{3}{@{}l}{\footnotesize\textit{Puissances}} \\
$P = U I\cos\varphi$ & Mono & W \\
$P = \sqrt{3}\,U_D I\cos\varphi$ & Tri & W \\
$S = \sqrt{3}\,U_D\,I$ & App. & VA \\
$I = P/(U\cos\varphi)$ & Mono & A \\
\midrule
\multicolumn{3}{@{}l}{\footnotesize\textit{Chute de tension}} \\
$\Delta U\% = \sqrt{3}IL(R_u c\varphi+X_u s\varphi)/U_D$ & Triphasé & \% \\
$\Delta U \leq 5\,\%$ & Dom. public & - \\
\midrule
\multicolumn{3}{@{}l}{\footnotesize\textit{Câble}} \\
$R_u = \rho/S$ & Rés. lin. & \si{\ohm/m} \\
\bottomrule
\end{tabular}

\columnbreak

\renewcommand{\arraystretch}{1.45}
\begin{tabular}{@{}p{4.5cm}ll@{}}
\toprule
\textbf{Formule} & \textbf{Usage} & \textbf{U.} \\
\midrule
\multicolumn{3}{@{}l}{\footnotesize\textit{DDR}} \\
Sél. : $\Delta I_{amont}/\Delta I_{aval} \geq 3$ & Select. & - \\
30\,mA - protection personnes & Obligat. & - \\
\midrule
\multicolumn{3}{@{}l}{\footnotesize\textit{Données typiques}} \\
$\rho_{Cu} = 1{,}7 \times 10^{-8}\,\Omega\!\cdot\!\mathrm{m}$ & Cuivre & - \\
$L = 850\,\mathrm{m}$ & Extracteur & m \\
$I_{cu} = 50\,\mathrm{kA}$ & NSX250N & - \\
$I_{rm} = 2860\,\mathrm{A}$ & MA220 & - \\
\midrule
\multicolumn{3}{@{}l}{\footnotesize\textit{Parafoudres}} \\
Type 1 - entrée installation & Impuls. & - \\
Type 2 - aval TGBT & Normal & - \\
TN-C : mode commun seulement & - & - \\
\bottomrule
\end{tabular}
\end{multicols}
\normalsize
\end{tcolorbox}

% -----------------------------------------------------------------------
%  3. QUESTIONS TYPES
% -----------------------------------------------------------------------

\begin{tcolorbox}[
  ficheSection=colElec,
  title={3. Questions types - extraites des sujets réels},
]

% -- QT1 -------------------------------------------------------------

\begin{tcolorbox}[questiontype,
  title={QT1 - Identification du SLT et choix du DDR \hfill\normalfont\footnotesize E2 disc. 2021 (Q19) + E2 disc. 2017 (B1.3)}
]
\textbf{Énoncé.}
L'installation est en régime TT. Justifier la présence d'un interrupteur différentiel et indiquer la sensibilité pour la protection des personnes.

\tcbline

\textbf{Corrigé.}

En régime \textbf{TT}, la prise de terre de l'installation (côté utilisateur) est indépendante de celle de la source. En cas de défaut d'isolement, le courant de défaut est limité par les résistances de prise de terre (typiquement quelques dizaines d'ohms) :
\[
  I_{defaut} = \frac{230}{R_{source} + R_{utilisateur}} \sim \text{quelques A} \ll I_{declenchement\,disjoncteur}
\]
Un disjoncteur seul ne peut pas détecter ce courant - le \textbf{DDR est obligatoire} selon NF C 15-100. Il mesure le courant différentiel $I_{diff} = \sum i_{conducteurs}$ et déclenche dès \textbf{30\,mA} pour la protection des personnes.

\textbf{Conclusion.} Sans DDR, le défaut d'isolement en régime TT resterait indéfiniment non détecté par la protection de tête : le DDR 30\,mA n'est pas une option mais une exigence normative directement liée à ce régime de neutre.
\end{tcolorbox}

\vspace{4pt}

% -- QT2 -------------------------------------------------------------

\begin{tcolorbox}[questiontype,
  title={QT2 - Calcul de la chute de tension et conformité NF C 15-100 \hfill\normalfont\footnotesize E1 2024 (B5-B9), Rapport jury 2024}
]
\textbf{Énoncé.}
Ligne triphasée de longueur $L = \SI{850}{\meter}$ alimentant un atelier de puissance
$P_{charge} = \SI{178,3}{\kilo\watt}$ compensé à $F_p = 1$, sous $U_r = \SI{400}{\volt}$
entre phases, courant de ligne $I_{tot} = \SI{258}{\ampere}$. Chaque conducteur est
modélisé par $R_L$ en série avec $L_L$, avec une résistance linéique
$R_u = \SI{20}{\milli\ohm\per\kilo\meter}$ et une inductance linéique
$L_u = \SI{1,2}{\milli\henry\per\kilo\meter}$ ($f = \SI{50}{\hertz}$).
Calculer $R_L$, $L_L$, puis $U_D$ et la chute de tension $\Delta U$ en proportion de $U_D$.
Conclure vis-à-vis de la norme NF C 15-100 (domaine public, «~autres usages~») et proposer
une solution si nécessaire.

\tcbline

\textbf{Corrigé.}

\textbf{Étape 1 — Paramètres de la ligne (B5).} La longueur vaut $L = \SI{0,850}{\kilo\meter}$ :
\[
  R_L = R_u \cdot L = 20\times10^{-3} \times 0{,}850 = \SI{17}{\milli\ohm}
  \qquad
  L_L = L_u \cdot L = 1{,}2\times10^{-3} \times 0{,}850 = \SI{1,02}{\milli\henry}
\]
L'énoncé fournit une \textbf{inductance} linéique et non une réactance : il faut la convertir
avec $\omega = 2\pi f = 2\pi \times 50 = \SI{314}{\radian\per\second}$ :
\[
  X_L = L_L\,\omega = 1{,}02\times10^{-3} \times 314 = \SI{0,320}{\ohm}
\]

\textbf{Étape 2 — Bilan de puissances vu du secondaire du transformateur (B7).} La charge étant
compensée ($Q_{charge} = 0$), seule la ligne consomme du réactif :
\[
  P_T = P_{charge} + 3R_L I_{tot}^2 = 178{,}3\times10^3 + 3\times0{,}017\times258^2
      \approx \SI{181,7}{\kilo\watt}
\]
\[
  Q_T = 3X_L I_{tot}^2 = 3\times0{,}320\times258^2 \approx \SI{64,0}{\kilo\volt\ampere}
  \qquad
  S_T = \sqrt{P_T^2 + Q_T^2} \approx \SI{192,6}{\kilo\volt\ampere}
\]

\textbf{Étape 3 — Tension au secondaire et chute relative (B8).}
\[
  U_D = \frac{S_T}{\sqrt{3}\,I_{tot}} = \frac{192{,}6\times10^3}{1{,}732\times258} \approx \SI{431}{\volt}
\]
\[
  \Delta U = U_D - U_r = 431 - 400 = \SI{31}{\volt}
  \qquad
  \frac{\Delta U}{U_D} = \frac{31}{431} \times 100 = \mathbf{7{,}2\,\%}
\]

\textbf{Conclusion (B9).} La norme NF C 15-100 impose, pour une installation alimentée par le
\textbf{réseau public} en catégorie «~\textbf{autres usages}~», une chute de tension
$\leq \SI{5}{\percent}$ (contre \SI{3}{\percent} pour l'éclairage). Ici $7{,}2\,\% > 5\,\%$ :
la ligne est \textbf{non conforme}. Solution : \textbf{augmenter la section} des conducteurs,
ce qui diminue simultanément $R_u = \rho/S$ et $L_u$, donc $\Delta U$, jusqu'à repasser sous le
seuil réglementaire.
\end{tcolorbox}

\vspace{4pt}

% -- QT3 -------------------------------------------------------------

\begin{tcolorbox}[questiontype,
  title={QT3 - Choix du disjoncteur associé à un variateur de fréquence \hfill\normalfont\footnotesize Rapport jury 2021 (Q17)}
]
\textbf{Énoncé.}
Choisir et justifier le disjoncteur à associer au variateur ATV630C11N4. Le courant de court-circuit réseau est $I_{cc} = 50\,\mathrm{kA}$.

\tcbline

\textbf{Corrigé.}

\textbf{Étape 1 — Référence.} NSX250\textbf{N} MA220.

\textbf{Étape 2 — Pouvoir de coupure.} Lettre \textbf{N} - $I_{cu} = 50\,\mathrm{kA}$ : adapté au courant de court-circuit du réseau.

\textbf{Étape 3 — Type de déclencheur.} \textbf{MA} (magnétique seul, seuil $I_{rm} = 2860\,\mathrm{A}$) : assure \emph{uniquement} la protection contre les courts-circuits. La protection thermique du moteur est prise en charge par le variateur (paramètre $I_{th}$ réglable de $0{,}2$ à $1{,}1\,I_n$).

\textbf{Conclusion.} Un disjoncteur magnéto-thermique créerait une double protection thermique en parallèle avec celle du variateur, avec un risque de décoordination (arrêt intempestif) : le choix MA est donc le seul conforme au cahier des charges d'un ensemble variateur.
\end{tcolorbox}

\end{tcolorbox}

% -----------------------------------------------------------------------
%  4. ATTENTES DU JURY
% -----------------------------------------------------------------------

\begin{tcolorbox}[
  jury,
  title={4. Attentes du jury - Erreurs récurrentes et points de vigilance},
]

\begin{enumerate}[label=\textbf{\arabic*.}, leftmargin=*, itemsep=5pt]

  \item \textbf{DDR obligatoire en TT - ne pas l'omettre.}
  En TT, le courant de défaut est trop faible pour déclencher un disjoncteur seul. Le DDR 30\,mA est obligatoire et la justification physique (résistances de prise de terre) doit être fournie.

  \item \textbf{TN-C : DDR interdit sur ce tronçon.}
  Le conducteur PEN confond N et PE - pas de courant différentiel mesurable - DDR inefficace et interdit sur le tronçon TN-C. Ne pas proposer le même équipement qu'en TN-S.

  \item \textbf{Partie réactive dans la chute de tension.}
  Si $X_u$ est fourni et $\sin\varphi \neq 0$, négliger $X_u \sin\varphi$ sous-estime $\Delta U$ de 10 à 20\,\%.

  \item \textbf{Conclure par rapport à la valeur normative.}
  Toujours comparer la valeur calculée à la limite (5\,\% pour autre usage, 3\,\% pour éclairage) et proposer une solution si non conforme.

  \item \textbf{Sélectivité différentielle : rapport $\geq 3$.}
  DDR 30\,mA aval - DDR amont $\geq 90\,\mathrm{mA}$ (soit 300\,mA en pratique). Ou utiliser un DDR de type S (temporisé) en amont.

  \item \textbf{Parafoudre TN-C.}
  En TN-C, seule la protection en mode commun est possible (N et PE confondus). Proposer uniquement mode commun ; la protection différentielle est impossible.

\end{enumerate}

\medskip
\begin{tcolorbox}[
  colback=colJuryFrame!12!colJury,
  colframe=colJuryFrame,
  left=4pt, right=4pt, top=2pt, bottom=2pt,
  fontupper=\small\itshape
]
Rapport jury 2024 : «~La chute de tension n'est pas respectée. La norme NFC 15-100 impose 5\,\% pour un autre usage (domaine public). Pour diminuer la chute de tension, il est possible de modifier les valeurs de $R_u$ et $L_u$ en augmentant la section du câble.~»\\
Rapport jury 2021 : «~Les candidats se sont limités à proposer une référence de composants parfois incomplète sans justification, sans indiquer clairement les critères de choix.~» — d'où l'exigence de justifier tout choix de disjoncteur par ses critères (pouvoir de coupure, type de déclencheur, réglage).\\
Rapport jury 2020 : «~Une critique à la fois qualitative et quantitative était attendue. Le jury déplore des réponses trop évasives sans lien avec les documents techniques.~»
\end{tcolorbox}

\end{tcolorbox}

% -----------------------------------------------------------------------
%  PIED - PRESENCE DANS LES SUJETS
% -----------------------------------------------------------------------

\begin{tcolorbox}[
  ficheSection=colPresence,
  title={Présence dans les sujets analysés (2013-2026)},
]
\begin{center}
\footnotesize\renewcommand{\arraystretch}{1.15}
\begin{tabular}{lccccccccccccc}
\toprule
\textbf{Année} & 13 & 14 & 15 & 16 & 17 & 18 & 19 & 20 & 21 & 22 & 23 & 24 & 25 \\
\midrule
E1 (commune)   & - & - & - & - & - & - & - & - & - & - & - & $\checkmark$ & - \\
E2 discipl. IE & - & - & - & - & $\checkmark$ & $\checkmark$ & - & $\checkmark$ & $\checkmark$ & - & - & $\checkmark$ & - \\
\bottomrule
\end{tabular}
\end{center}
\smallskip
\footnotesize\textbf{Sous-thèmes fréquents :}
SLT TT et DDR 30\,mA (5/13) - chute de tension NF C 15-100 (4/13) - disjoncteur MA amont variateur (4/13) - sélectivité différentielle (3/13) - parafoudres TT/TN/IT (2/13)

\medskip
\textbf{Ordres de grandeur (repères rapides).}
\begin{itemize}[leftmargin=1.3em]
  \item Réseau BT FR : tension phase-neutre typique \SI{230}{\volt}, tension phase-phase typique \SI{400}{\volt}, fréquence \SI{50}{\hertz}.
  \item Courants d'étude fréquents en sujet : de quelques \si{\ampere} à quelques dizaines de \si{\ampere}.
  \item En conversion d'énergie, un rendement attendu reste généralement entre 0{,}8 et 0{,}98 selon la technologie.
\end{itemize}

\medskip
\textbf{Synthèse express (à refaire sans support).}
\begin{enumerate}[leftmargin=1.4em]
  \item Poser les hypothèses du modèle et les conventions de signe.
  \item Écrire 2 à 4 relations clés avec unités.
  \item Réaliser un calcul numérique complet et vérifier l'ordre de grandeur.
  \item Justifier le résultat par un argument technique de niveau CAPET.
\end{enumerate}

\end{tcolorbox}

\CAPETENDFILE
